\chapter{Agoraphobia}
\index{Agoraphobia}

\section*{Definition and Overview}
Agoraphobia is an anxiety disorder in which a person fears and avoids situations where escape may be difficult or where help may not be available if panic-like symptoms or other embarrassing symptoms occur. Typical feared situations include crowds, public transport, supermarkets, open spaces, queues, and being outside the home alone.

Many people with agoraphobia also have panic disorder, and the two conditions frequently begin together: a first panic attack in a particular setting leads to fear of that setting, which then spreads to similar situations. In severe cases the person becomes effectively housebound, unable to leave home without a trusted companion. Agoraphobia is a genuine, treatable condition and not a sign of weakness or poor character.

\section*{Epidemiology}
Agoraphobia affects roughly 1 to 2 percent of adults at some point in their lives and is about twice as common in women as in men. Onset is most often in the late teens to early 30s, although it can begin at any age. Agoraphobia frequently develops after a first panic attack, and many affected people also experience other anxiety disorders or depression.

Without treatment the condition often runs a chronic course, but it can improve substantially with appropriate therapy. Many people delay seeking help for years because of shame or because they interpret their symptoms as a physical illness.

\section*{Aetiology and Pathophysiology}
Agoraphobia usually develops through a combination of biological, psychological, and environmental factors.

\begin{itemize}
    \item \textbf{Panic attacks:} a frightening first panic attack, followed by fear of having another in a place where escape is difficult
    \item \textbf{Avoidance learning:} avoiding the feared situation brings short-term relief, which powerfully reinforces the avoidance
    \item \textbf{Temperament:} a tendency to be anxious, or to be highly alert to bodily sensations such as a racing heart or dizziness
    \item \textbf{Genetics and family history:} anxiety disorders run in families, suggesting a heritable component
    \item \textbf{Stressful events:} loss, illness, or major life changes can trigger onset
    \item \textbf{Physical conditions:} disorders causing dizziness, faintness, or palpitations can fuel the fear
    \item \textbf{Childhood experiences:} overprotective parenting and learning anxious responses from family members
\end{itemize}

\section*{Clinical Features}
\begin{itemize}
    \item Fear of specific situations such as buses, trains, supermarkets, queues, or being far from home
    \item Avoiding these situations, needing a companion, or enduring them only with intense distress
    \item Panic symptoms when confronting the situation: pounding heart, sweating, shaking, shortness of breath, dizziness, or nausea
    \item Fear of losing control, fainting, or embarrassing themselves in public
    \item Anticipatory anxiety -- worrying about a feared situation well before it occurs
    \item In severe cases, being largely confined to the home
\end{itemize}

\section*{Differential Diagnosis}
\begin{itemize}
    \item Panic disorder, in which panic attacks occur but avoidance of situations is not prominent
    \item Social anxiety disorder, in which the fear is of being judged by others rather than of escape
    \item Specific phobias, such as fear of lifts or flying
    \item Depression, which can cause withdrawal and difficulty leaving the house
    \item Medical conditions causing dizziness, faintness, or palpitations, such as heart disease, thyroid problems, or inner-ear disorders
    \item Substance-related anxiety, such as from caffeine, stimulants, or alcohol withdrawal
\end{itemize}

\section*{When to Seek Medical Care}
Seek help if avoidance is restricting work, study, or social life, if panic attacks are frequent, or if the person is becoming housebound. Agoraphobia is very treatable, and earlier help tends to work better.

\textbf{Contact your local emergency services for:}
\begin{itemize}
    \item Panic symptoms accompanied by chest pain, difficulty breathing, or fainting, to exclude a physical emergency
    \item Thoughts of suicide or self-harm
\end{itemize}

\section*{Diagnosis and Investigations}
\begin{itemize}
    \item \textbf{Clinical interview:} discussion of feared situations, avoidance behaviour, panic symptoms, and the impact on daily life
    \item \textbf{Questionnaires:} validated scales for anxiety and agoraphobia severity to gauge the diagnosis and monitor progress
    \item \textbf{Physical examination:} to look for medical causes of the physical symptoms
    \item \textbf{Blood tests:} thyroid function and other tests to exclude physical conditions
    \item \textbf{ECG:} sometimes used to check the heart when palpitations or chest discomfort are prominent
\end{itemize}

\section*{Management}
\begin{itemize}
    \item \textbf{Cognitive behaviour therapy (CBT):} the main treatment, combining gradual exposure to feared situations with challenging frightening thoughts
    \item \textbf{Relaxation and breathing techniques:} to manage the physical symptoms of anxiety
    \item \textbf{Antidepressant medication:} selective serotonin reuptake inhibitors may be prescribed for moderate to severe agoraphobia, usually alongside therapy
    \item \textbf{Support:} family involvement, and structured practice with a trusted companion where needed
    \item \textbf{Self-help and online CBT:} structured programmes can help with mild to moderate symptoms
    \item \textbf{Treating underlying problems:} such as depression or physical conditions contributing to symptoms
\end{itemize}

\section*{Complications}
\begin{itemize}
    \item Becoming housebound and dependent on others
    \item Depression and additional anxiety disorders
    \item Alcohol or sedative misuse as a form of self-medication
    \item Social isolation and relationship difficulties
    \item Loss of work, study, or driving ability
    \item Reduced quality of life and physical deconditioning from inactivity
\end{itemize}

\section*{Prognosis and Outcomes}
With treatment, most people with agoraphobia improve significantly. Cognitive behaviour therapy with gradual exposure is effective in many cases, and combining it with medication can help in more severe cases. Untreated, the disorder often persists for many years with a fluctuating course. Relapses can occur during stressful periods, but learning and practising coping skills reduces the risk and helps people maintain the ground they have gained.

\section*{Prevention and Self-Care}
\begin{itemize}
    \item \textbf{Early help:} seek help early if panic attacks or avoidance begin
    \item \textbf{Gradual exposure:} practise facing feared situations in small, manageable steps rather than avoiding them
    \item \textbf{Calming techniques:} use calm breathing and grounding techniques during anxious moments
    \item \textbf{Substances:} limit caffeine, and avoid alcohol and stimulants that can trigger panic symptoms
    \item \textbf{Lifestyle:} keep up exercise, sleep, and routine to reduce overall anxiety
    \item \textbf{Connections:} maintain social contact and continue everyday activities as much as possible
    \item \textbf{Support:} if you care for someone with agoraphobia, encourage small steps and celebrate progress
\end{itemize}

\section*{Sources and Further Reading}
The following sources provide reliable, up-to-date information on agoraphobia.

\begin{itemize}
    \item NHS. \textit{Conditions: agoraphobia.} https://www.nhs.uk/conditions/
    \item Mayo Clinic. \textit{Agoraphobia: symptoms, causes, and treatment.} https://www.mayoclinic.org
    \item MedlinePlus. \textit{Agoraphobia.} https://medlineplus.gov
    \item MSD Manual. \textit{Agoraphobia -- overview and treatment.} https://www.msdmanuals.com
    \item National Institute of Mental Health. \textit{Anxiety disorders including panic and agoraphobia.} https://www.nimh.nih.gov
    \item World Health Organization. \textit{Mental health and anxiety disorders resources.} https://www.who.int
\end{itemize}
